%=====================================================================
\section{Shapes, sizes and run sets}\label{sec:ssetup}
\paragraph{Why these shapes.} Every cavity here is simply connected with one wall, so the imposed winding $k$ fixes the same
total charge as in a circle. Changing the shape at fixed $k$ (and fixed $\rho$) separates topology from geometry (H1).
\begin{itemize}
  \item \textbf{Cassini ovals} $d_1d_2=b^2$ with foci at $(\pm a,0)$ (Sec.~\ref{sec:md-cassini}): a one-parameter family from the
        circle ($a=0$) through convex ovals ($b/a\ge\sqrt2$) to a concave neck ($1<b/a<\sqrt2$) and the lemniscate ($b=a$).
        Half-length $\sqrt{a^2+b^2}$, half-height $b^2/(2a)$ for $b/a<\sqrt2$, waist $2\sqrt{b^2-a^2}$. Aspect ratio = length/height.
        Winding imposed uniformly in arclength (eMode 3). Families are labelled by their half-length (H34, H48, H64 cells $=1.53$,
        2.16, $2.88\,\la$) and $b/a$.
  \item \textbf{Ellipses} with semi-axes $A=\varepsilon B$ (along $x$) and $B$, area 2500 ($\sqrt{A/\pi}=1.27\,\la$), arclength-uniform
        winding (eMode 2); $\varepsilon=1.2$, 1.5, 2.0. A second convex family with a different curvature profile (no flattened sides).
\end{itemize}

\begin{table}[H]\centering\small
\caption{Oval and ellipse run sets ($\zeta=0.07$ unless passive). Steps: warm-up + production. All in \texttt{runs/cassini/}.}\label{tab:ssets}
\begin{tabularx}{\textwidth}{llXcl}\toprule
set & cavities & shapes / parameters & seeds & steps\\\midrule
T1, T2 & validation & passive $k=3$: $a=0$, $b/a=1.6$, 1.2; eMode 0 and 3 & 1 & short\\
P1 & H48, H64 & $b/a=1.6$, 1.3, 1.1; $k=6$ & 1 & 10k + 40k\\
P2 & H48, H64 pinched & $b/a=1.01$--1.05, $k=6$; passive $b/a=1.02$, 1.01 & 1 & 10k + 40k\\
P3 & H34 small ovals & $b/a=2.5$, 2.0, 1.6, 1.3, $k=6$ & 3 & 10k + 100k\\
P3 & H48 pinched, high $k$ & $b/a=1.02$, $k=12$, 15, 18, 20; passive $k=18$, 20 & 1 & 10k + 100k\\
P4 & lobe-shaped oval & H30.4, $b/a=1.8$ (area and aspect of one lobe of H64 $b/a=1.01$), $k=3$ & 3 & 10k + 100k\\
P5 & passive references & H34 $b/a=2.5$--1.3 and the three ellipses, $k=6$ (5 Oct) & 1 & 10k + 30k\\
W1 & H48 size series & $b/a=2.5$, 2.0, 1.6, 1.3, $k=6$ & 3 & 10k + 100k\\
W2 & H34 $k$ scan & $b/a=2.0$, 1.6; $k=2$, 4, 8, 10, 12 ($k=6$ = P3) & 3 & 10k + 100k\\
W3 & ellipses & $\varepsilon=1.2$, 1.5, 2.0, $k=6$, eMode 2 & 2 & 10k + 100k\\
W4 & H64 seeds & $b/a=1.3$, 1.1: seeds 2, 3 (seed 1 = P1) & 2 & 10k + 40k\\\bottomrule
\end{tabularx}
\end{table}

\paragraph{Matching ovals with circles.} An oval has no single radius, so the comparison is made in every size measure at once,
at matched linear defect density $\rho=2k/P$ (Table~\ref{tab:match}). The test is whether an oval lies inside the region where
circles keep the layer (0.75--$2\,\la$, Sec.~\ref{sec:sweep}) in all of them.
\begin{table}[H]\centering\small
\caption{Ovals and circles: density and candidate size measures (units of $\la$).}\label{tab:match}
\begin{adjustbox}{max width=\textwidth}
\begin{tabular}{lccccccc}\toprule
cavity & $k$ & aspect & $\rho=2k/P$ & $\sqrt{A/\pi}$ & half-length & inscribed $r$ & $2A/P$\\\midrule
circle $R=1.43\,\la$ & 6 & 1 & 0.060 & 1.43 & 1.43 & 1.43 & 1.43\\
circle $R=1.07\,\la$ & 3 & 1 & 0.040 & 1.07 & 1.07 & 1.07 & 1.07\\
oval H34 $b/a=2.5$ & 6 & 1.17 & 0.060 & 1.42 & 1.53 & 1.30 & 1.41\\
oval H34 $b/a=1.6$ & 6 & 1.51 & 0.066 & 1.27 & 1.53 & 1.01 & 1.23\\
oval H34 $b/a=1.3$ & 6 & 1.94 & 0.069 & 1.15 & 1.53 & 0.79 & 1.07\\
oval H48 $b/a=1.6$ & 6 & 1.51 & 0.046 & 1.80 & 2.16 & 1.43 & 1.74\\
lobe-shaped oval & 3 & 1.38 & 0.036 & 1.18 & 1.37 & 1.00 & 1.16\\\bottomrule
\end{tabular}
\end{adjustbox}
\end{table}
The smallest measure, the inscribed radius, is 0.79--1.01\,$\la$ for the small ovals that lose the layer; circles of 0.72 and
0.96\,$\la$ at the same density keep a clean layer (Sec.~\ref{sec:sweep}). Every elongated oval therefore sits inside the circles'
layer region in every size measure.
